\section{Outlook: three grand challenges}

The next stage of embodied medicine will not be defined by a single more capable robot or foundation model. Progress depends on solving three linked challenges: systems that can discover and validate interventions, patient models that remain trustworthy over time, and closed-loop therapies that can learn without sacrificing safety.

\paragraph{From automated execution to autonomous biomedical discovery.}
Self-driving laboratories already combine robotic experimentation, machine learning and closed-loop optimization to navigate chemical and materials spaces \cite{Hase2019NextGenerationExperimentation,Burger2020MobileRoboticChemist,MacLeod2020SelfDrivingLaboratory,Angello2022ClosedLoopOptimization,Tom2024SelfDrivingLaboratories}. Translating this paradigm to medicine is harder because the objective is not merely yield or material performance: relevant outcomes may be delayed, heterogeneous, ethically constrained and partly unobservable. A clinically useful autonomous discovery system must connect mechanistic hypotheses, biological models, assay choice, uncertainty-aware experimental design and human review. It must also recognize when an experiment is invalid, when a surrogate endpoint is misleading, and when evidence is insufficient to proceed. Medical-agent benchmarks and autonomous decision systems are early steps toward evaluating sequential reasoning and tool use, but virtual records and curated tasks do not yet reproduce the physical, regulatory and interpersonal constraints of care \cite{Jiang2025MedAgentBench,Ferber2025OncologyAgent,Ferber2026AutonomousMedicalAIAgents,Tu2025ConversationalDiagnosticAI}. The grand challenge is therefore auditable co-discovery: machines propose and execute bounded experiments, while clinicians and scientists retain authority over aims, stopping rules and translation.

\paragraph{From static digital twins to longitudinal patient models.}
Personal omics and longitudinal profiling established that health is a dynamic trajectory across molecular, physiological and behavioral scales \cite{Chen2012PersonalOmics,SchuesslerFiorenzaRose2019LongitudinalBigData,Zhou2019LongitudinalMultiOmics}. Wearables extend this view into daily life and can predict laboratory measurements or detect departures from an individual's baseline \cite{Dunn2021WearablePredictions,Mishra2020SmartwatchCOVID}. Yet a digital twin is useful only if it updates under distribution shift, represents uncertainty and distinguishes disease progression from changes in devices, behavior or care. Multimodal biomedical AI can integrate heterogeneous signals, but missingness and incompatible timescales remain part of the phenomenon rather than a nuisance to be imputed away \cite{Acosta2022MultimodalBiomedicalAI,Katsoulakis2024DigitalTwins}. The required patient model should be causal enough to support counterfactual intervention questions, modest enough to abstain when unidentifiable, and portable enough to survive transitions between home, clinic and hospital. Its validation target is not resemblance to historical data but prospective usefulness across time and settings.

\paragraph{From fixed controllers to adaptive therapeutic partnerships.}
Closed-loop insulin delivery and responsive neuromodulation show that sensing and intervention can be coupled safely when the target, actuator and safeguards are sharply specified \cite{Brown2019ClosedLoopDiabetes,Scangos2021ClosedLoopNeuromodulation,FDA2023PhysiologicClosedLoop}. Emerging platforms combine biochemical sensing with drug delivery or support programmable recording and stimulation during natural behavior \cite{Lee2016GrapheneMicroneedles,Topalovic2023ClosedLoopWearable}. The harder problem is continual adaptation: physiology changes, sensors drift, preferences evolve and multiple therapies interact. An embodied partner must learn within a safety envelope, separate personalization from uncontrolled experimentation, and preserve meaningful patient and clinician control. This calls for hierarchical architectures in which adaptive policies operate beneath verified constraints, updates are logged and reversible, and escalation is triggered by uncertainty or discordant evidence. Learning health systems provide the institutional counterpart, turning carefully governed outcomes into model improvement rather than silent deployment drift \cite{McDonald2024LearningHealthSystem}.

Together these challenges define a shift from isolated intelligent devices to accountable learning ecosystems. The decisive benchmark will not be autonomy in the abstract, but whether the system can generate better evidence, maintain a faithful and contestable model of the individual, and intervene with measurable benefit under real-world constraints.

